\begin{pastexam}[label={px:ntu-108-comm-dsab-27}]{臺大電信所 108 學年度（B）第 27 題}
（單選；每題 2.5 分，答錯不倒扣）

Also note that in this exam, the degree of a node $x$ in a tree indicates the number of children
of $x$.

What is the maximum possible degree of a root in a Fibonacci heap with 256 nodes?

A.~8 \qquad B.~9 \qquad C.~10 \qquad D.~11 \qquad E.~16
\end{pastexam}

\begin{pastsoln}
答案是 D。

\textbf{上界}：費氏堆積裡階數 $k$ 的節點，子樹至少有 $F_{k+2}$ 個節點（資結書第~8.6~節）。 理由：把
$x$ 的子節點依被連結的先後排成 $y_1, \ldots, y_k$。 $y_i$ 被連結到 $x$ 底下時，$x$ 已經有 $y_1,
\ldots, y_{i-1}$，階數至少 $i - 1$，而連結只發生在同階的兩個根之間，所以 $y_i$ 當時的階數也至少
$i - 1$；之後 $y_i$ 最多失去一個子節點（失去第二個就會被剪下），所以現在階數至少 $i - 2$。 設 $s_k$
是階數 $k$ 的節點子樹最少的節點數，得到 $s_k \ge 2 + \sum_{i=2}^{k} s_{i-2}$，由此 $s_k \ge F_{k+2}$。

費氏數：$F_{12} = 144$、$F_{13} = 233$、$F_{14} = 377$。 一共只有 $256$ 個節點，階數 $12$ 要
$F_{14} = 377 > 256$ 個，不可能；階數 $11$ 要 $F_{13} = 233 \le 256$ 個。 也可以寫成
$k \le \lfloor \log_\phi 256 \rfloor = \lfloor 11.52 \rfloor = 11$。

\textbf{做得到}：連續插入 $2049$ 個元素再取出最小值，整理後剩一棵 $B_{11}$（$2048 = 2^{11}$），根 $r$
的階數 $11$，子節點是 $B_{10}, \ldots, B_0$ 的根。 接著修剪：對 $r$ 的每個子節點 $c$（原本是 $B_j$ 的根），把它最大的那棵
子樹 $B_{j-1}$ 的根的鍵減到 $-\infty$，整棵剪到根串列，再一個一個取出最小值把剪下來的節點全部刪掉；
然後對 $c$ 剩下的子節點遞迴做同樣的事。 每個非根節點只失去一個子節點，只會被標記，不會引發連鎖
剪下；剪下來的樹最多 $2^9$ 個節點，整理時只會彼此連結，階數到不了 $11$，不會和 $r$ 連在一起。 修剪後 $B_j$ 剩 $F_{j+1}$ 個節點，$r$ 的子樹剩
$1 + (F_1 + F_2 + \cdots + F_{11}) = 1 + (F_{13} - 1) = 233$ 個，$r$ 的階數仍是 $11$。 最後再插入
$23$ 個元素（插入不整理），湊滿 $256$ 個。 所以最大可能是 $11$。

A 的 $8 = \lg 256$ 是二項堆積的答案，費氏堆積因為會剪下子樹，樹可以比較「稀疏」，同樣節點數能撐出
更大的階數；B、C 不是最大；E 做不到。
\end{pastsoln}
